\documentclass[12pt]{article}
\usepackage{amssymb}
\usepackage{amsmath}
\usepackage{amsthm}
\usepackage{array}
\usepackage{graphicx}
  \DeclareGraphicsRule{*}{mps}{*}{}
%\usepackage{bbm}
\usepackage{marvosym}
\usepackage{enumerate}
\usepackage{txfonts}
\usepackage{paralist}
\usepackage{pdfpages}
\usepackage{multicol}
\usepackage{pgfplots}
\pgfplotsset{compat=1.5}
\everymath{\displaystyle}

\usepackage{hyperref}
\hypersetup{
    colorlinks=true,
    urlcolor=blue,
}

%\usepackage{showframe}

\pagestyle{empty}

\usepackage{fullpage}
%% \addtolength{\hoffset}{-1.5cm}
%% \addtolength{\textwidth}{3cm}

%% \addtolength{\voffset}{-1.5cm}
%% \addtolength{\textheight}{3cm}

\setdefaultleftmargin{0pt}{}{}{}{}{}

\newcommand{\vv}[2]{#1}

\hyphenpenalty=5000
\tolerance=1000
\setlength{\parindent}{0pt}

\begin{document}

\thispagestyle{empty}
\begin{center}
  \textbf{ Math 202 -- Slice, approximate, integrate: volumes of solids of revolution
  --- Fall 2024
    }
\end{center}

For each of the regions listed below:

\begin{itemize}
    \item Draw a graph of all the relevant curves. (Using Desmos or something is a good choice.)
    \item Find the intersection points of the curves. (Using Desmos or something is a good choice.)
    \item Shade in the region that's ``bounded by'' these curves.
    \item Draw the axis of revolution and sketch the solid that results. (Do ur best; this is not an art class.)
    \item Draw a representative slice perpendicular to the axis of revolution; revolve the slice.
    \item Label the radius (radii?) of the slice and the thickness of the slice.
    \item Use your labels to write an equation for $\Delta V$.
\end{itemize}

\vfill
\href{https://activecalculus.org/single/sec-6-2-volume.html#act-6-2-1}{Activity 6.2.2} \hrulefill

\begin{enumerate}[a.]
\item The region bounded by the $x$-axis, $y=\sqrt{x}$, and $x=4$; revolve about the $x$-axis.
\item The region bounded by the $y$-axis, $y=\sqrt{x}$, and $y=2$; revolve about the $x$-axis.
\item The region bounded by $y=\sqrt{x}$ and $y=x^3$; revolve about the $x$-axis.
\item The region bounded by $y=2x^2+1$ and $y=x^2 + 4$; revolve about the $x$-axis.

%% I omitted part (e) because it's the first part of the next activity.

\vfill

\hspace{-1.5em} \href{https://activecalculus.org/single/sec-6-2-volume.html#act-6-2-2}{Activity 6.2.3} \hrulefill

%(a)
\item The region bounded by the $y$-axis, $y=\sqrt{x}$, and $y=2$; revolve about the $y$-axis.
%(b)
\item The region bounded by the $x$-axis, $y=\sqrt{x}$, and $x=4$; revolve about the $y$-axis.
%(c)
\item The region in the first quadrant bounded by $y=2x$ and $y=x^3$; revolve about the $x$-axis.
%(d)
\item The region in the first quadrant bounded by $y=2x$ and $y=x^3$; revolve about the $y$-axis.
%(e)
\item The region bounded by $x=(y-1)^2$ and $y=x-1$; revolve about the $y$-axis.

\vfill

\hspace{-1.5em} \href{https://activecalculus.org/single/sec-6-2-volume.html#act-6-2-3}{Activity 6.2.4} \hrulefill
%(a)
\item The region in the first quadrant bounded by $y=2x$ and $y=x^3$; revolve about the line $y=-2$.
%(b)
\item The same region, but revolve about the line $y=4$.
%(c)
\item The same region, but revolve about the line $x=-1$.
%(d)
\item The same region, but revolve about the line $x=5$.

\end{enumerate}

\end{document}
